%! Unit 7: The Singular Value Decomposition — Summative Assessment — Practice Test --- Study Copy (KEY)
\documentclass[10pt]{article}
\usepackage{linalg-article}
\usepackage{linalg-key}   % pulls in -boxes; do NOT also load -boxes
\newcommand{\xx}{\mathbf{x}}
\newcommand{\vv}{\mathbf{v}}
\newcommand{\uu}{\mathbf{u}}
\newcommand{\T}{^{\mathsf T}}

% Part-divider strip
\newcommand{\parthead}[1]{%
  \vspace{6pt}\noindent
  \begin{headlinebox}{burgundy}{\color{white}\bfseries #1}\end{headlinebox}%
  \vspace{4pt}}

\begin{document}

\pageheader{Unit 7: The Singular Value Decomposition}{Practice Test --- Study Copy --- Answer Key}
\namedateperiod

\vspace{0.10in}
\begin{remindbox}
\small
\textbf{This is a practice test.} It mirrors the real Unit 7 test in length, format,
and the ideas it covers, but the numbers and contexts differ. Work every problem
with no notes, then check your answers against the key.
\end{remindbox}

% ======================================================================
\parthead{Part A --- Vocabulary (8 pts)}
Write the letter of the best description next to each term.

\vspace{4pt}
\noindent
\begin{minipage}[t]{0.40\textwidth}
\begin{enumerate}[leftmargin=2.2em, itemsep=7pt, label=\arabic*.]
  \item Singular value \ans{C}
  \item Singular vectors \ans{E}
  \item Singular Value Decomposition \ans{B}
  \item Outer product \ans{A}
  \item Rank-one matrix \ans{F}
  \item Orthogonal matrix \ans{D}
  \item Principal component \ans{H}
  \item Energy of a matrix \ans{G}
\end{enumerate}
\end{minipage}%
\hfill
\begin{minipage}[t]{0.57\textwidth}
\small
\begin{description}[leftmargin=1.4em, itemsep=3pt]
  \item[(A)] A column times a row, $\uu\vv\T$ --- a whole matrix built from two vectors.
  \item[(B)] The factorization $A=U\Sigma V\T$ (rotate--stretch--rotate) that \emph{every} matrix has.
  \item[(C)] A number $\sigma\ge0$: how far $A$ stretches along a singular direction (a half-width of the ellipse).
  \item[(D)] A square matrix with perpendicular unit columns, so $Q\T Q=I$ and $Q^{-1}=Q\T$.
  \item[(E)] The perpendicular unit inputs $\vv_i$ and outputs $\uu_i$ with $A\vv_i=\sigma_i\uu_i$.
  \item[(F)] A matrix whose every column is a multiple of one vector (like $\uu\vv\T$).
  \item[(G)] The total $\sigma_1^2+\sigma_2^2+\cdots$, equal to the sum of the squares of all entries of $A$.
  \item[(H)] A top singular direction of centered data --- the direction of greatest spread.
\end{description}
\end{minipage}

% ======================================================================
\parthead{Part B --- Multiple Choice (12 pts)}
Circle the best answer.

\begin{enumerate}[leftmargin=1.9em, itemsep=9pt]
  \item The singular values $\sigma$ of a matrix are
    \hfill \textbf{(A) always $\ge0$} \ans{$\leftarrow$} \quad (B) sometimes negative \quad (C) the eigenvalues of $A$ \quad (D) always $1$
  \item To find the singular values of $A$, use the eigenvalues of
    \hfill (A) $A$ \quad (B) $A+A\T$ \quad \textbf{(C) $A\T A$} \ans{$\leftarrow$} \quad (D) $A^{-1}$
  \item A singular value $\sigma_i$ and the eigenvalue $\lambda_i$ of $A\T A$ are related by
    \hfill (A) $\sigma=\lambda$ \quad \textbf{(B) $\sigma=\sqrt{\lambda}$} \ans{$\leftarrow$} \quad (C) $\sigma=\lambda^2$ \quad (D) $\sigma=1/\lambda$
  \item In $A=U\Sigma V\T$, all the stretching is carried by
    \hfill (A) $U$ \quad (B) $V$ \quad \textbf{(C) $\Sigma$} \ans{$\leftarrow$} \quad (D) $V\T$
  \item An orthogonal matrix $Q$ satisfies
    \hfill (A) $Q^2=I$ \quad \textbf{(B) $Q\T Q=I$, so $Q^{-1}=Q\T$} \ans{$\leftarrow$} \quad (C) $\det Q=0$ \quad (D) $Q=Q\T$ always
  \item To compress an image with the SVD, you keep the rank-one layers with the
    \hfill (A) smallest $\sigma$ \quad \textbf{(B) largest $\sigma$} \ans{$\leftarrow$} \quad (C) largest $\uu$ \quad (D) negative $\sigma$
\end{enumerate}

\begin{teachernote}
Item 1: singular values are always $\ge0$ (they are $\sqrt{\lambda}$, lengths of the ellipse's half-axes). Item 2:
$\sigma$'s come from the eigenvalues of the symmetric $A\T A$. Item 3: $\sigma=\sqrt{\lambda}$. Item 4: $U$ and $V\T$ are
rotations (they only turn); $\Sigma$ is the only factor that stretches. Item 5: orthogonal means $Q\T Q=I$, so the
inverse is free, $Q^{-1}=Q\T$. Item 6: compression keeps the biggest-$\sigma$ layers (they hold the most energy).
\end{teachernote}

% ======================================================================
\parthead{Part C --- Short Answer \& Computation (35 pts)}
Show your work.

\begin{enumerate}[leftmargin=1.9em, itemsep=10pt]
  \item Singular values from geometry.
        \ans{\textbf{(a)} A diagonal matrix stretches by its diagonal entries, so the singular values are
        $\sigma_1=4,\ \sigma_2=3$ (largest first). The unit circle becomes an ellipse with half-widths $4$ and $3$.
        \textbf{(b)} The columns of $S=\begin{bsmallmatrix}0&2\\1&0\end{bsmallmatrix}$ are $\begin{bsmallmatrix}0\\1\end{bsmallmatrix}$
        (length $1$) and $\begin{bsmallmatrix}2\\0\end{bsmallmatrix}$ (length $2$), and they are perpendicular, so
        $\sigma_1=2,\ \sigma_2=1$.}
  \item Singular values of $A=\begin{bsmallmatrix}2&2\\1&-2\end{bsmallmatrix}$ via $A\T A$.
        \ans{$A\T A=\begin{bsmallmatrix}5&2\\2&8\end{bsmallmatrix}$ (entries are the column dot products).
        $\det(A\T A-\lambda I)=(5-\lambda)(8-\lambda)-4=\lambda^2-13\lambda+36=(\lambda-9)(\lambda-4)=0$, so
        $\lambda=9,4$. Therefore $\sigma_1=\sqrt9=3$ and $\sigma_2=\sqrt4=2$.}
  \item Output vector $\uu_1=\dfrac{A\vv_1}{\sigma_1}$.
        \ans{$A\vv_1=\dfrac{1}{\sqrt5}\begin{bsmallmatrix}2&2\\1&-2\end{bsmallmatrix}\begin{bsmallmatrix}1\\2\end{bsmallmatrix}
        =\dfrac{1}{\sqrt5}\begin{bsmallmatrix}6\\-3\end{bsmallmatrix}$. Divide by $\sigma_1=3$:
        $\uu_1=\dfrac{1}{3\sqrt5}\begin{bsmallmatrix}6\\-3\end{bsmallmatrix}=\dfrac{1}{\sqrt5}\begin{bsmallmatrix}2\\-1\end{bsmallmatrix}$.
        \textbf{Unit check:} $|\uu_1|^2=\tfrac{1}{5}(2^2+(-1)^2)=\tfrac{5}{5}=1$. \checkmark}
  \item Outer product and rank one.
        \ans{\textbf{(a)} $\uu\vv\T=\begin{bsmallmatrix}1\\3\end{bsmallmatrix}\begin{bsmallmatrix}2&1\end{bsmallmatrix}
        =\begin{bsmallmatrix}2&1\\6&3\end{bsmallmatrix}$. \textbf{(b)} Both columns are multiples of
        $\uu=\begin{bsmallmatrix}1\\3\end{bsmallmatrix}$ (namely $2\uu$ and $1\uu$), so every column points along one
        line --- the matrix has \textbf{rank one}.}
  \item Rank-one approximation and energy of $A=\begin{bsmallmatrix}4&2\\2&4\end{bsmallmatrix}$.
        \ans{\textbf{(a)} $A_1=\sigma_1\uu_1\vv_1\T=6\cdot\tfrac12\begin{bsmallmatrix}1\\1\end{bsmallmatrix}\begin{bsmallmatrix}1&1\end{bsmallmatrix}
        =\begin{bsmallmatrix}3&3\\3&3\end{bsmallmatrix}$. \textbf{(b)} Total energy $=\sigma_1^2+\sigma_2^2=36+4=40$
        (equal to $4^2+2^2+2^2+4^2=40$). $A_1$ keeps $\dfrac{36}{40}=90\%$ of the energy.}
  \item Orthogonal matrix $Q=\dfrac15\begin{bsmallmatrix}3&-4\\4&3\end{bsmallmatrix}$.
        \ans{\textbf{(a)} The columns are $\tfrac15\begin{bsmallmatrix}3\\4\end{bsmallmatrix}$ and
        $\tfrac15\begin{bsmallmatrix}-4\\3\end{bsmallmatrix}$: each has length $\tfrac15\sqrt{9+16}=1$, and their dot
        product is $\tfrac{1}{25}(-12+12)=0$. So $Q\T Q=I$ ($1$'s on the diagonal, $0$'s off). \textbf{(b)}
        $Q^{-1}=Q\T=\dfrac15\begin{bsmallmatrix}3&4\\-4&3\end{bsmallmatrix}$ (the inverse is free). \textbf{(c)}
        $Q\xx=\tfrac15\begin{bsmallmatrix}3&-4\\4&3\end{bsmallmatrix}\begin{bsmallmatrix}5\\0\end{bsmallmatrix}
        =\begin{bsmallmatrix}3\\4\end{bsmallmatrix}$, which has length $\sqrt{9+16}=5$ --- \textbf{the same length as
        $\xx$}, because an orthogonal matrix is a rigid motion (rotation/reflection) that preserves length.}
  \item PCA of $(5,7),(7,5),(3,1),(1,3)$.
        \ans{\textbf{(a)} Mean $=(4,4)$; centered points $(1,3),(3,1),(-1,-3),(-3,-1)$. \textbf{(b)}
        $A\T A=\begin{bsmallmatrix}20&12\\12&20\end{bsmallmatrix}$ (column dot products of the centered data).
        \textbf{(c)} $\lambda=20\pm12=32,8$; the top eigenvector is $\mathrm{PC}_1=\tfrac{1}{\sqrt2}\begin{bsmallmatrix}1\\1\end{bsmallmatrix}$.
        \textbf{(d)} $\mathrm{PC}_1$ captures $\dfrac{32}{32+8}=\dfrac{32}{40}=80\%$ of the spread.}
\end{enumerate}

% ======================================================================
\parthead{Part D --- Extended Response (10 pts)}
Explain your reasoning in full sentences.

\begin{enumerate}[leftmargin=1.9em, itemsep=12pt]
  \item Why every matrix has singular values.
        \ans{\textbf{(a)} $A\T A$ is symmetric because $(A\T A)\T=A\T(A\T)\T=A\T A$ --- it equals its own transpose no
        matter what shape $A$ is. By the spectral theorem a symmetric matrix has \emph{real} eigenvalues and
        \emph{perpendicular} eigenvectors; and each eigenvalue satisfies $\lambda=\dfrac{|A\vv|^2}{|\vv|^2}\ge0$
        (a length squared over a length squared), so $\lambda\ge0$. \textbf{(b)} Since every $\lambda\ge0$, its square
        root $\sigma=\sqrt{\lambda}$ is a real number $\ge0$. These $\sigma$'s are the singular values --- and because
        $A\T A$ exists and is symmetric for \emph{any} matrix, every matrix has them (unlike eigenvectors, which can be
        complex or too few).}
  \item Rotate--stretch--rotate and orthogonal matrices.
        \ans{\textbf{(a)} Reading $A=U\Sigma V\T$ right to left: $V\T$ \emph{rotates} the input frame onto the axes,
        $\Sigma$ \emph{stretches} along those axes by the singular values $\sigma_1,\sigma_2,\dots$, and $U$
        \emph{rotates} the result to the output frame. $U$ and $V\T$ are orthogonal, so they only turn (no resizing);
        all the stretching is packed into the diagonal $\Sigma$. \textbf{(b)} $Q\T Q=I$ says $Q\T$ times $Q$ gives the
        identity, which is exactly what it means for $Q\T$ to be the inverse: $Q^{-1}=Q\T$, free of any computation.
        Length is preserved because $|Q\xx|^2=(Q\xx)\T(Q\xx)=\xx\T Q\T Q\xx=\xx\T\xx=|\xx|^2$, so $|Q\xx|=|\xx|$.}
\end{enumerate}

\begin{teachernote}
\textbf{Scoring Part D (5 pts each).} D1: 3 pts (a) $A\T A$ symmetric for any $A$ $\Rightarrow$ real, perpendicular
eigenvectors and $\lambda\ge0$ (via $\lambda=|A\vv|^2/|\vv|^2$); 2 pts (b) $\sigma=\sqrt\lambda\ge0$ real, so every
matrix has singular values. D2: 3 pts (a) $V\T$/$U$ rotate, $\Sigma$ stretches, all stretching in $\Sigma$; 2 pts (b)
$Q\T Q=I\Rightarrow Q^{-1}=Q\T$, and $|Q\xx|^2=\xx\T Q\T Q\xx=|\xx|^2$. Accept equivalent wording throughout.
\end{teachernote}

\end{document}
